\sect{Note: a quadratic lower bound for sign-coherent circuits (preliminary)}\label{app:coherent}

\begin{keybox}
\textbf{Status.} This appendix records a recent result together with
everything needed to reproduce it. The proofs below are complete and have
been checked by an independent referee and by computer on small cases. The
constants are not optimised, and the result is not yet formalized in Lean.
We plan to tighten and integrate it later.
\end{keybox}

\smalltitle{The class.}
A bottom gate is \emph{sign-coherent} if, inside every address block, its
weights are either all $\ge0$ or all $\le0$. Its table weights and the top
gate are arbitrary. Every gate that sees a block only through $|a_i|$ is
sign-coherent, and so is every gate that is monotone in the address bits.
All gates of the circuit in Theorem~\ref{thm:exact} are of this kind.

\begin{theorem}\label{thm:coherent}
Let $C$ be a depth-two threshold circuit for $A$ whose $S$ bottom gates are
all sign-coherent. Put $D_0=\bigl(N-k-1-\log_2(S+1)\bigr)/\bigl(k+\log_2(S+1)\bigr)$. Then
\[
S\ \ge\ \frac{m^k}{m^k-(m-1)^k}\,D_0\ \ge\ \frac mk\,D_0 .
\]
Consequently $S\ge\min\{N^2,\ m(N-3k-1)/(3k^2)\}$, and in the balanced case
\[
S\ \ge\ \frac{N(N-3k-1)}{3k^3}\ =\ \Omega\!\left(\frac{n^2}{\log^3 n}\right).
\]
\end{theorem}

So within this class the circuit of Theorem~\ref{thm:exact} is optimal up to a
polylogarithmic factor. A circuit with substantially fewer gates must contain
gates whose weights have mixed signs inside a block.

\smalltitle{Plan.}
We walk through all unit cubes of the grid of block weights. On each cube,
the circuit reveals the whole table, so many gates must be non-constant there
(Lemma~\ref{lem:entropy}). A sign-coherent gate is non-constant on only a
$k/m$ fraction of the cubes (Lemma~\ref{lem:diag}). Comparing the two counts
gives the theorem.

\smalltitle{The cube restriction.}
Fix an order of the positions inside every block. For a vector
$s\in\{0,\dots,m\}^k$ let $a(s)$ be the address part in which block $i$ has
ones exactly at its first $s_i$ positions. A \emph{cube} is a vector
$c\in\{0,\dots,m-1\}^k$; there are $m^k$ cubes. For $z\in\bits^k$ the input
$(x,a(c+z))$ differs from $(x,a(c))$ exactly in the live positions
$c_i+1$ of the blocks $i$ with $z_i=1$, and its address is
$(c\bmod2)\oplus z$. Hence
\[
z\mapsto A\bigl(x,a(c+z)\bigr)=x_{(c\bmod 2)\oplus z}
\]
determines the whole table $x$. A bottom gate \emph{splits} the cube $c$ for
the table $x$ if it is non-constant in $z$ on the inputs $(x,a(c+z))$. Let
$d_c(x)$ be the number of bottom gates that split $c$ for $x$.

\begin{lemma}[Many gates split each cube]\label{lem:entropy}
Let $C$ be any depth-two threshold circuit for $A$ with $S$ bottom gates, and
let $x$ be uniformly random. Then $\mathbb E_x[d_c(x)]\ge D_0$ for every cube
$c$.
\end{lemma}
\begin{proof}
Fix $c$. On the inputs $(x,a(c+z))$, bottom gate $j$ computes
$\ind{\lambda_j\cdot z\ge\tau_j}$, where $\lambda_j\in\mathbb R^k$ lists its
weights on the live positions and only the number $\tau_j$ depends on $x$.
Likewise the top gate computes
$\ind{\sum_jv_jg_j+\mu\cdot z\ge\kappa}$, where $v_j$ are its weights on the
bottom gates and $\mu$ its weights on the live positions; the gates that
do not split $c$ are constant and are absorbed, together with the table and
the fixed address bits, into $\kappa$. So only $\tau_j$ and $\kappa$ depend
on $x$.

Therefore the function $z\mapsto x_{(c\bmod2)\oplus z}$, and with it $x$, is
determined by: the set $G$ of splitting gates; for each $j\in G$ the
non-constant function $\ind{\lambda_j\cdot z\ge\tau_j}$, one of at most $2^k-1$
(only the position of $\tau_j$ among the $\le2^k$ values $\lambda_j\cdot z$
matters); and the position of $\kappa$, one of at most $2^k+1$. With
$D=|G|=d_c(x)$, entropy gives
\[
N=H(x)\le H(D)+\mathbb E\log_2\Bigl[\tbinom SD(2^k-1)^D(2^k+1)\Bigr]
\le\log_2(S+1)+\mathbb E[D]\,(\log_2S+k)+k+1,
\]
which rearranges to the claim.
\end{proof}

\begin{lemma}[A coherent gate splits few cubes]\label{lem:diag}
For every table $x$, a sign-coherent gate splits at most $m^k-(m-1)^k$ cubes.
\end{lemma}
\begin{proof}
Let $\varepsilon_i=+1$ if the gate's weights in block $i$ are $\ge0$ and
$\varepsilon_i=-1$ otherwise. Its linear form on $(x,a(s))$ is a constant plus
$\sum_iP_i(s_i)$, where $P_i(s_i)$ is the sum of its weights on the first
$s_i$ positions of block $i$. So the form $F(s)$ does not decrease when
$s_i$ moves by one step in direction $\varepsilon_i$.

On the cube $c$ the form is smallest at $c^-=c+(\mathbf1-\varepsilon)/2$ and
largest at $c^+=c+(\mathbf1+\varepsilon)/2$. The gate with threshold $\theta$
splits $c$ exactly when $F(c^-)<\theta\le F(c^+)$.

Group the cubes into \emph{lines} $c,\,c+\varepsilon,\,c+2\varepsilon,\dots$.
For consecutive cubes on a line, the largest vertex of one is the smallest
vertex of the next, so the intervals $(F(c^-),F(c^+)]$ along a line are
disjoint, and $\theta$ lies in at most one of them. Each line starts at a cube
$c$ for which $c-\varepsilon$ is not a cube, and exactly $(m-1)^k$ cubes $c$
have $c-\varepsilon$ a cube. So there are $m^k-(m-1)^k$ lines.
\end{proof}

\begin{proof}[Proof of Theorem~\ref{thm:coherent}]
Sum Lemma~\ref{lem:entropy} over all cubes and compare with
Lemma~\ref{lem:diag}: $m^kD_0\le\sum_c\mathbb E[d_c]\le S\,(m^k-(m-1)^k)$.
Then use $1-(1-1/m)^k\le k/m$. If $S<N^2$ then $\log_2(S+1)\le2k$ and
$D_0\ge(N-3k-1)/(3k)$. In the balanced case $m=N/k$; this requires $k$ to be
a power of two.
\end{proof}

\begin{remark}
(1) For gates that see each block only through $|a_i|$, a slab argument with
Ball's cube-slicing bound improves the denominator to $6\sqrt2\,k^{5/2}$ for
$k\ge8$. (2) Lemma~\ref{lem:diag} is tight for all $k,m$. The fact that a
unate function changes at most once along a chain compatible with its
orientation is classical (Impagliazzo, Paturi and Saks~\cite{IPS},
Lemma~2.1); the cube restriction and the quadratic bound appear to be new,
but our literature search was limited. (3) Theorem~\ref{thm:coherent} is
asymptotic. It exceeds $n^{3/2}$ only from about $k=37$, and the method
(Lemma~\ref{lem:entropy} plus a uniform bound as in Lemma~\ref{lem:diag})
cannot give more than $O(n^2/\log n)$. (4) Lemma~\ref{lem:entropy} holds for
all circuits. For general gates, the survival-type arguments of
Section~\ref{sec:lower} cannot exceed $O(n^{3/2})$: under every
distribution of restrictions that leaves one live bit per block, a gate with
random $\pm1$ weights is non-constant with probability at least
$1/(2\sqrt{2km})$. Gates with mixed signs inside a block are exactly where
the gap between $n^{3/2}$ and $n^2$ lies.
\end{remark}

\smalltitle{Two address blocks.}
For $k=2$ the same ideas give matching linear bounds.

\begin{proposition}\label{prop:k2}
Let $k=2$. (a) Every circuit for $A$ with sign-coherent bottom gates has at
least $m-1$ bottom gates. (b) There is a circuit with $4\lceil m/2\rceil+m$
bottom gates, each seeing the blocks only through $s_i=|a_i|$.
\end{proposition}
\begin{proof}[Proof sketch of (a)]
Take the table with $x_c=c_1\oplus c_2$. On every cube (here a unit square)
the restricted circuit computes XOR or its negation. If every bottom gate were
constant or a literal there, the top gate would compute XOR as a threshold
function, which is impossible. So on every square some gate cuts off exactly
one corner. After reflecting $s_i\mapsto m-s_i$ where needed, the region where
a coherent gate is $1$ is an up-set of the grid $\{0,\dots,m\}^2$. Its
boundary is a monotone lattice path, and a corner is cut exactly where the
path turns. On the $4(m-1)$ squares at the border of the grid such a path
turns at most four times. So at least $m-1$ gates are needed.
\end{proof}
\begin{proof}[Construction for (b)]
Write $b=b_1+2b_2$ for the table index and $J=\lceil m/2\rceil$. The bottom
gates are, for $1\le j\le J$,
\[
\ind{s_1+x_0\ge2j},\ \ind{s_1+x_2\ge2j},\ \ind{s_2+x_0\ge2j},\ \ind{s_2+x_1\ge2j}
\]
with top weights $-4,-2,-4,-2$, and, for $1\le j\le m$, $\ind{s_1+s_2+x_0\ge2j}$
with top weight $-4$. The top gate outputs $\ind{F\ge0}$ with
\[
F=\sum_j v_jg_j+5s_1+5s_2+8x_0+2x_1+\tfrac32x_2+x_3-\tfrac{13}2 .
\]
For $0\le t\le T$ we have $t-2\sum_{j\le\lceil T/2\rceil}\ind{t+y\ge2j}=(-1)^y(t\bmod2)$
for every bit $y$, so $F$ depends on $s$ only through $b$. This leaves a
finite check, and we verified $\ind{F\ge0}=x_b$ for all inputs with exact
arithmetic for $m\le40$.
\end{proof}

The earlier general constructions give about $12m$ gates for $k=2$. Exact
solver data for gates that see blocks only through $|a_i|$ give minima $3,3$
for $m=1,2$ and lower bounds $5$ for $m=3,4,5$ and $6$ for $m=6,8,10$.

\smalltitle{Reproducing.}
Lemma~\ref{lem:diag} can be checked by brute force: enumerate all cubes and
compute $F(c^\pm)$ for random sign-coherent weight vectors, including zero
weights and ties, and random orders. Equality is attained by a gate with one
large weight on the last position of each block. For the bound of (a), the
Ring Lemma can be checked over all up-sets of $\{0,\dots,m\}^2$ for
$m\le11$. The construction of (b) is checked as described. All three checks
were run independently twice.
